%==========================================================
%  CHAPITRE 2 — Analyse et spécification des besoins
%  CRM Intelligent pour Centres d'Appels
%==========================================================

\chapter{Analyse et spécification des besoins}
\chapterplan{%
\chapterplanitem{1}{Introduction}{sec:chap2_introduction}
\chapterplanitem{2}{Spécification des besoins}{sec:specification_besoins}
\chapterplanitem{3}{Diagramme de cas d'utilisation global}{sec:diagramme_cas_utilisation}
\chapterplanitem{4}{Diagramme de classes}{sec:diagramme_classes}
\chapterplanitem{5}{Product Backlog}{sec:product_backlog}
\chapterplanitem{6}{Répartition des sprints}{sec:repartition_sprints}
\chapterplanitem{7}{Conclusion}{sec:chap2_conclusion}
}

%----------------------------------------------------------
\section{Introduction}
\phantomsection
\label{sec:chap2_introduction}
%----------------------------------------------------------

La phase d'analyse et de spécification des besoins constitue une étape fondamentale dans la
réussite de tout projet logiciel. Elle permet de définir précisément les attentes des différents
utilisateurs ainsi que les contraintes techniques et fonctionnelles indispensables à la conception
d'un système robuste et adapté.

Ce chapitre présente l'identification des besoins de la plateforme CRM intelligente, la
modélisation des fonctionnalités principales à travers plusieurs diagrammes UML, ainsi que la
planification du projet selon la méthodologie Scrum. \cite{ScrumGuide}

%==========================================================
\section{Spécification des besoins}
\label{sec:specification_besoins}
%==========================================================

%---------------------------------------------------------
\subsection{Identification des acteurs}
%---------------------------------------------------------

La plateforme CRM met en interaction cinq acteurs principaux dont les rôles et les
responsabilités sont définis selon une hiérarchie de droits progressifs :

\begin{itemize}[label=\textbullet, itemsep=5pt]

    \item \textbf{Agent :}
    utilisateur principal du centre d'appels. Il peut consulter et gérer ses leads, planifier
    des rendez-vous, passer des appels, consulter son scoring IA et son planning de présence.
    Ce rôle constitue la base de la hiérarchie opérationnelle.

    \item \textbf{Quality Manager :}
    responsable de l'évaluation de la qualité des appels. Il hérite des droits de consultation
    de l'Agent et peut en plus évaluer manuellement les appels, consulter les scores IA,
    générer des rapports de qualité et configurer les règles d'évaluation.

    \item \textbf{Admin RH :}
    responsable de la gestion des ressources humaines. Il gère les agents (création,
    modification, désactivation), les plannings de présence, les pauses et le calcul des
    salaires selon les règles définies.

    \item \textbf{Supervisor :}
    superviseur des opérations du centre d'appels. Il consulte les tableaux de bord de
    performance, suit les campagnes, analyse les métriques de conversion et peut accéder
    à l'ensemble des données opérationnelles de son périmètre.

    \item \textbf{Super Admin :}
    administrateur global de la plateforme. Il dispose de privilèges d'administration sur
    l'ensemble du système : gestion des utilisateurs et de leurs rôles, configuration des
    permissions, paramétrage des règles de scoring IA, des poids d'évaluation et des
    configurations système.

\end{itemize}

%---------------------------------------------------------
\subsection{Spécification des besoins fonctionnels par acteur}
%---------------------------------------------------------

\subsubsection{L'Agent}

\begin{itemize}[label=--, itemsep=3pt]
    \item \textbf{Gestion des accès :}
    \begin{itemize}[label=--, itemsep=2pt]
        \item S'authentifier et se déconnecter de manière sécurisée.
        \item Accéder à son espace de travail personnel.
    \end{itemize}

    \item \textbf{Gestion des leads :}
    \begin{itemize}[label=--, itemsep=2pt]
        \item Consulter la liste des leads qui lui sont assignés.
        \item Filtrer les leads par statut, campagne ou date.
        \item Ajouter un nouveau lead avec ses informations de contact.
        \item Modifier les informations d'un lead existant.
        \item Qualifier un lead (HORS_CIBLE, RDV1, RDV2, RAPPEL, REFUS, etc.).
        \item Planifier un follow-up sur un lead.
    \end{itemize}

    \item \textbf{Gestion des rendez-vous :}
    \begin{itemize}[label=--, itemsep=2pt]
        \item Consulter son agenda de rendez-vous.
        \item Créer un rendez-vous pour un lead.
        \item Modifier ou annuler un rendez-vous.
    \end{itemize}

    \item \textbf{Gestion des appels :}
    \begin{itemize}[label=--, itemsep=2pt]
        \item Consulter l'historique de ses appels.
        \item Visualiser la transcription de ses appels.
        \item Consulter son score IA par appel.
    \end{itemize}

    \item \textbf{Consultation IA :}
    \begin{itemize}[label=--, itemsep=2pt]
        \item Visualiser son score IA global et par critère.
        \item Consulter son évolution dans le temps.
    \end{itemize}

    \item \textbf{Présence :}
    \begin{itemize}[label=--, itemsep=2pt]
        \item Pointer son arrivée et son départ.
        \item Déclarer une pause.
    \end{itemize}
\end{itemize}

\subsubsection{Le Quality Manager}

\begin{itemize}[label=--, itemsep=3pt]
    \item \textbf{Évaluation qualité :}
    \begin{itemize}[label=--, itemsep=2pt]
        \item Consulter la liste des appels à évaluer.
        \item Effectuer une évaluation manuelle d'un appel avec notes et commentaires.
        \item Comparer son évaluation avec le score IA automatique.
        \item Consulter l'historique des évaluations.
    \end{itemize}

    \item \textbf{Scoring IA :}
    \begin{itemize}[label=--, itemsep=2pt]
        \item Visualiser les détails du scoring IA par appel.
        \item Analyser les tendances de qualité par agent et par équipe.
    \end{itemize}

    \item \textbf{Alertes :}
    \begin{itemize}[label=--, itemsep=2pt]
        \item Configurer des règles d'alerte (score minimum, inactivité, taux de conversion).
        \item Consulter les alertes déclenchées.
    \end{itemize}
\end{itemize}

\subsubsection{L'Admin RH}

\begin{itemize}[label=--, itemsep=3pt]
    \item \textbf{Gestion des agents :}
    \begin{itemize}[label=--, itemsep=2pt]
        \item Créer, modifier et désactiver des comptes agents.
        \item Consulter la liste des agents.
    \end{itemize}

    \item \textbf{Gestion de présence :}
    \begin{itemize}[label=--, itemsep=2pt]
        \item Consulter les pointages de présence.
        \item Visualiser les statistiques de présence par agent et par période.
    \end{itemize}

    \item \textbf{Gestion de la paie :}
    \begin{itemize}[label=--, itemsep=2pt]
        \item Définir les règles de calcul des salaires.
        \item Consulter les salaires calculés.
        \item Exporter les rapports de paie.
    \end{itemize}
\end{itemize}

\subsubsection{Le Supervisor}

\begin{itemize}[label=--, itemsep=3pt]
    \item \textbf{Tableaux de bord :}
    \begin{itemize}[label=--, itemsep=2pt]
        \item Consulter les KPIs globaux du centre d'appels.
        \item Visualiser les performances par agent, par équipe et par campagne.
        \item Analyser les taux de conversion et les tendances.
    \end{itemize}

    \item \textbf{Campagnes :}
    \begin{itemize}[label=--, itemsep=2pt]
        \item Créer et gérer des campagnes.
        \item Suivre les performances des campagnes.
    \end{itemize}

    \item \textbf{Rapports :}
    \begin{itemize}[label=--, itemsep=2pt]
        \item Générer des exports PDF des rapports de performance.
        \item Exporter les données d'analyse.
    \end{itemize}
\end{itemize}

\subsubsection{Le Super Admin}

\begin{itemize}[label=--, itemsep=3pt]
    \item \textbf{Configuration globale :}
    \begin{itemize}[label=--, itemsep=2pt]
        \item Gérer les utilisateurs et leurs rôles.
        \item Configurer les permissions granulaires.
        \item Paramétrer les poids des critères de scoring IA.
        \item Configurer les règles d'alerte globales.
        \item Gérer les paramètres de transcription et de sentiment.
    \end{itemize}

    \item \textbf{Supervision :}
    \begin{itemize}[label=--, itemsep=2pt]
        \item Accéder à l'ensemble des données de la plateforme.
        \item Surveiller l'état du système.
    \end{itemize}
\end{itemize}

%---------------------------------------------------------
\subsection{Besoins non fonctionnels}
%---------------------------------------------------------

\begin{itemize}[label=\textbullet, itemsep=5pt]

    \item \textbf{Sécurité :}
    \begin{itemize}[itemsep=2pt]
        \item Authentification JWT avec tokens d'accès et de refresh.
        \item Contrôle d'accès basé sur les permissions (RBAC).
        \item Chiffrement des communications (TLS 1.3).
        \item Hachage des mots de passe.
        \item Conformité RGPD avec anonymisation des données.
    \end{itemize}

    \item \textbf{Performance :}
    \begin{itemize}[itemsep=2pt]
        \item Temps de réponse API inférieur à 500ms pour les requêtes standards.
        \item Scoring IA avec timeout configurable (120s max).
        \item Transcription audio avec modèle optimisé (base, français).
    \end{itemize}

    \item \textbf{Disponibilité :}
    \begin{itemize}[itemsep=2pt]
        \item Architecture Dockerisée avec healthchecks.
        \item Base de données PostgreSQL avec persistance.
        \item Déploiement via Render ou auto-hébergé.
    \end{itemize}

    \item \textbf{Maintenabilité :}
    \begin{itemize}[itemsep=2pt]
        \item Architecture Clean Architecture avec séparation des couches.
        \item Repository Pattern et Unit of Work.
        \item Documentation Swagger/OpenAPI des endpoints.
        \item Tests unitaires, d'intégration et E2E.
    \end{itemize}

    \item \textbf{Extensibilité :}
    \begin{itemize}[itemsep=2pt]
        \item API REST versionnée.
        \item Ajout simplifié de nouveaux critères de scoring.
        \item Configuration des poids d'évaluation via fichier JSON.
    \end{itemize}

\end{itemize}

%==========================================================
\section{Diagramme de cas d'utilisation global}
\label{sec:diagramme_cas_utilisation}
%==========================================================

Le diagramme de cas d'utilisation global illustre les interactions entre les cinq acteurs et
les fonctionnalités principales de la plateforme CRM.

\begin{figure}[H]
    \centering
    \fbox{\parbox{0.95\linewidth}{
    \centering\vspace{0.4cm}
    \textit{[Diagramme de cas d'utilisation global de la plateforme CRM]}\\[0.3cm]
    Ce diagramme présente les cinq acteurs (Super Admin, Admin RH, Supervisor, Quality Manager
    et Agent) et leurs cas d'utilisation principaux~: authentification, gestion des leads, des
    rendez-vous, des appels, du scoring IA, de la qualité, de la présence, de la paie, des
    campagnes et de l'administration système.
    \vspace{0.4cm}
    }}
    \caption{Diagramme de cas d'utilisation global de la plateforme CRM}
    \label{fig:use_case_crm}
\end{figure}

%==========================================================
\section{Diagramme de classes}
\label{sec:diagramme_classes}
%==========================================================

Le diagramme de classes modélise les 26 entités de la plateforme CRM ainsi que leurs
relations. Il constitue la base de la conception orientée objet en traduisant les besoins
fonctionnels en un modèle conceptuel cohérent. Les entités sont organisées en cinq
paquetages selon leur domaine métier.

%% ──────────────────────────────────────────────────────────
\subsubsection{Paquetage 1 : Authentification et Sécurité}
%% ──────────────────────────────────────────────────────────

Ce paquetage regroupe les entités responsables de la gestion des identités, des rôles
et des permissions granulaires.

\begin{description}[itemsep=4pt]
    \item[\textbf{User}] : entité centrale du système. Elle stocke les informations
    d'authentification (Username, Password hashé avec BCrypt, Email), le profil (Name)
    et le rôle (\texttt{UserRole} : Admin, Agent, Qualite). Elle est reliée à la plupart
    des autres entités du système~:
    \begin{itemize}
        \item \textbf{1 ---* Attendance} : un utilisateur a plusieurs pointages
        \item \textbf{1 ---* CrmAppointment} : un agent a plusieurs rendez-vous clients
        \item \textbf{1 ---* Salary} : un agent a un salaire par mois
        \item \textbf{1 ---* ManualEvaluation (Agent)} : un agent reçoit des évaluations
        \item \textbf{1 ---* ManualEvaluation (Evaluator)} : un évaluateur donne des évaluations
        \item \textbf{1 ---* Message (Sender)} : un utilisateur envoie des messages
        \item \textbf{1 ---* Message (Receiver)} : un utilisateur reçoit des messages
    \end{itemize}

    \item[\textbf{Permission}] : représente une permission granulaire (ex~: \texttt{Leads.Create},
    \texttt{Scoring.Configure}). Chaque permission appartient à un module (Leads, Agenda,
    Scoring, etc.) et possède une description. Elle est liée aux rôles via l'entité de
    jointure \texttt{RolePermission}.

    \item[\textbf{RolePermission}] : entité de jointure entre un rôle (Admin, Agent, Qualite)
    et une permission. Un même rôle peut avoir plusieurs permissions (ex~: le rôle Admin
    possède toutes les permissions), et une même permission peut être attribuée à
    plusieurs rôles.
\end{description}

\textbf{Relations du paquetage~:}
\begin{itemize}
    \item Permission \textbf{1 ---* } RolePermission (une permission est assignée à plusieurs rôles)
    \item RolePermission \textbf{* ---* } User (via le champ Role)
\end{itemize}

%% ──────────────────────────────────────────────────────────
\subsubsection{Paquetage 2 : Gestion Commerciale}
%% ──────────────────────────────────────────────────────────

Ce paquetage couvre le coeur métier du CRM~: la gestion des prospects, des appels,
des rendez-vous et des campagnes.

\begin{description}[itemsep=4pt]
    \item[\textbf{Lead}] : représente un prospect commercial. Il contient ses coordonnées
    (CompanyName, ContactName, Phone, Email, Address, PostalCode), son statut (Status)
    et l'identifiant de l'agent auquel il est affecté (\texttt{AgentId}). Le statut
    évolue selon la qualification (nouveau, RDV1, RDV2, REFUS, RAPPEL, HORS\_CIBLE).

    \item[\textbf{Call}] : enregistrement d'un appel téléphonique. C'est l'entité la plus
    riche du modèle avec plus de 50 attributs~:
    \begin{itemize}
        \item \textbf{Transcription} : texte brut (Transcription), textes séparés
        agent/client (AgentText, ClientText), transcription labellisée (LabeledTranscript)
        \item \textbf{Scores IA} : ScorePercentage (score global), 8 scores par critère
        (ScoreEcoute, ScorePersuasion, ScoreEmpathie, ScoreArgumentation, ScoreRefus,
        ScoreVente, AgentPoliteness, etc.)
        \item \textbf{Analyse} : Sentiment (positif/négatif/neutre), SentimentScore,
        Summary, Keywords, CustomerIntent
        \item \textbf{Métadonnées} : CallDuration, CallDate, AgentTalkRatio, AgentSeconds,
        InactivityDetected, InactivityDuration, RefusalReason, QualificationCoherence
    \end{itemize}

    \item[\textbf{CrmAppointment}] : rendez-vous client planifié par un agent. Il contient
    les informations détaillées du projet (ProjectType, Chauffage, Toiture, Isolation,
    Consommation, CreditScore, Revenus), le statut financier (FinancingStatus, SituationBancaire),
    la qualité (QualityScore) et le statut (pending/confirmed/cancelled).

    \item[\textbf{Appointment}] : rendez-vous détecté automatiquement par l'IA depuis
    un appel. Stocke la date détectée (DetectedDate), le score de confiance
    (ConfidenceScore) et le statut (detected/confirmed/rejected).

    \item[\textbf{Campaign}] : campagne marketing avec son nom (Name) et le nom
    de l'entreprise ciblée (CompanyName).

    \item[\textbf{Followup}] : suivi de relance d'un lead. Contient le nom de l'agent
    (AgentName), la date de rendez-vous (AppointmentDate), le statut et le nombre de
    relances (RelanceCount).

    \item[\textbf{LeadFolder}] : dossier de classement pour organiser les leads
    (Name, Campaign).

    \item[\textbf{Qualification}] : table de référence des codes de qualification
    (Code : HORS\_CIBLE, RDV1, RDV2, REFUS, RAPPEL) avec leur libellé (Label) et
    les mots-clés attendus (ExpectedKeywords) pour la validation IA.

    \item[\textbf{Supplier}] : fournisseur associé aux leads ou campagnes.
\end{description}

\textbf{Relations du paquetage~:}
\begin{itemize}
    \item Lead \textbf{* --- 1} User (un lead est affecté à un agent)
    \item Lead \textbf{1 --- *} Call (un lead peut avoir plusieurs appels)
    \item Lead \textbf{1 --- *} Followup (un lead peut avoir plusieurs relances)
    \item Call \textbf{1 --- *} Appointment (un appel peut générer plusieurs RDV détectés)
    \item Lead \textbf{* --- 1} Campaign (un lead appartient à une campagne)
    \item Lead \textbf{* --- 1} LeadFolder (un lead peut être classé dans un dossier)
\end{itemize}

%% ──────────────────────────────────────────────────────────
\subsubsection{Paquetage 3 : Présence et Paie}
%% ──────────────────────────────────────────────────────────

Ce paquetage gère les ressources humaines~: pointage des agents, pauses et calcul
automatisé des salaires.

\begin{description}[itemsep=4pt]
    \item[\textbf{Attendance}] : pointage journalier d'un agent. Enregistre la date
    (Date), l'heure d'arrivée (ClockIn), l'heure de départ (ClockOut) et le statut
    (active/completed). Lié à User et contient une collection de pauses.

    \item[\textbf{AttendanceBreak}] : pause effectuée pendant une journée de travail.
    Stocke le début (BreakStart), la fin (BreakEnd) et le type (Type~: déjeuner,
    café, technique). La durée de la pause est déduite du temps de travail total.

    \item[\textbf{Salary}] : salaire mensuel calculé pour un agent. Contient le
    détail complet~: salaire de base (BaseSalary), nombre de RDV (RdvCount),
    nombre de poses (PoseCount), nombre de refus (RefusCount), taux qualité
    (QualityRate), primes (RdvBonus, PoseBonus, QualityBonus, InstallationBonus),
    pénalités (Penalties) et salaire total (TotalSalary). Le statut de paiement
    (PaymentStatus) peut être pending/paid/cancelled.

    \item[\textbf{SalaryRule}] : règle de calcul de salaire configurable. Définit
    le type de règle (RuleType~: rdv\_bonus, pose\_bonus, quality\_bonus,
    penalty), le montant (Amount), le rôle ciblé (Role) et l'état actif/inactif
    (IsActive).
\end{description}

\textbf{Relations du paquetage~:}
\begin{itemize}
    \item User \textbf{1 --- *} Attendance (un agent a un pointage par jour)
    \item Attendance \textbf{1 --- *} AttendanceBreak (un pointage contient plusieurs pauses)
    \item User \textbf{1 --- *} Salary (un agent a un salaire par mois)
\end{itemize}

%% ──────────────────────────────────────────────────────────
\subsubsection{Paquetage 4 : Qualité et Intelligence Artificielle}
%% ──────────────────────────────────────────────────────────

Ce paquetage regroupe les entités liées à l'évaluation de la qualité des appels,
au scoring IA et aux alertes.

\begin{description}[itemsep=4pt]
    \item[\textbf{ManualEvaluation}] : évaluation manuelle de la qualité d'un appel
    réalisée par un Quality Manager ou Admin. Contient l'identifiant de l'agent
    évalué (AgentId), de l'évaluateur (EvaluatorId), la note globale (GlobalScore),
    la décision (Decision), les commentaires (Commentaires) et les scores détaillés
    au format JSON (ScoresJson). Un agent peut recevoir plusieurs évaluations,
    et un évaluateur peut en donner plusieurs.

    \item[\textbf{AlertRule}] : règle d'alerte qualité configurable. Définit le type
    (RuleType~: score\_min, inactivite, conversion), le seuil (ThresholdValue),
    l'email de notification (NotificationEmail) et l'état actif (IsActive).

    \item[\textbf{AlertHistory}] : historique des alertes déclenchées. Stocke le nom
    de l'agent concerné (AgentName), le type d'alerte (AlertType), la sévérité
    (Severity~: warning/critical), le message, le seuil et la valeur réelle
    (ActualValue).

    \item[\textbf{AiEligibilityLog}] : journal des vérifications d'éligibilité
    effectuées par l'IA. Stocke l'identifiant de l'agent (AgentId), les données
    client (ClientData) et le résultat (Result).
\end{description}

\textbf{Relations du paquetage~:}
\begin{itemize}
    \item User \textbf{1 --- *} ManualEvaluation (en tant qu'agent évalué)
    \item User \textbf{1 --- *} ManualEvaluation (en tant qu'évaluateur)
    \item User \textbf{1 --- *} AiEligibilityLog (log des vérifications par agent)
\end{itemize}

%% ──────────────────────────────────────────────────────────
\subsubsection{Paquetage 5 : Communication et Système}
%% ──────────────────────────────────────────────────────────

Ce paquetage contient les entités de messagerie interne et les données auxiliaires.

\begin{description}[itemsep=4pt]
    \item[\textbf{Message}] : message de la messagerie interne temps réel (WebSocket).
    Contient l'expéditeur (SenderId), le destinataire (ReceiverId), le contenu
    (Content), le flag d'urgence (IsUrgent), l'état de lecture (IsRead) et les
    horodatages (CreatedAt, ReadAt).

    \item[\textbf{AgentSavedData}] : données spécifiques sauvegardées par un agent
    (préférences, configuration personnelle). Lié à User.

    \item[\textbf{Agent}] : statistiques agrégées d'un agent utilisées pour les
    rapports (FirstCall, LastCall, TotalCalls).

    \item[\textbf{SourceFile}] : fichier source importé dans la plateforme
    (fichiers audio, documents).

    \item[\textbf{Log}] : journalisation système pour le suivi des actions.
\end{description}

\textbf{Relations du paquetage~:}
\begin{itemize}
    \item User \textbf{1 --- *} Message (Sender) : un utilisateur envoie plusieurs messages
    \item User \textbf{1 --- *} Message (Receiver) : un utilisateur reçoit plusieurs messages
    \item User \textbf{1 --- *} AgentSavedData : un agent a des données sauvegardées
\end{itemize}

%% ──────────────────────────────────────────────────────────
\subsection{Synthèse des relations}
%% ──────────────────────────────────────────────────────────

\begin{figure}[H]
    \centering
    \fbox{\parbox{0.95\linewidth}{
    \centering\vspace{0.3cm}
    {\large\textbf{Diagramme de classes global de la plateforme CRM}}\\[0.2cm]
    \textit{26 entités réparties en 5 paquetages}\\[0.3cm]
    \small
    \renewcommand{\arraystretch}{1.4}
    \begin{tabular}{|c|c|c|}
    \hline
    \rowcolor{tekupBlue!20}
    \textbf{Paquetage} & \textbf{Entité} & \textbf{Attributs principaux} \\
    \hline
    \multirow{3}{2.8cm}{\centering Authentification \\ \& Sécurité} & \textbf{User} & Id, Username, Password, Name, Role (Admin|Agent|Qualite), Email \\
    \cline{2-3}
    & \textbf{Permission} & Id, Name, Module, Description \\
    \cline{2-3}
    & \textbf{RolePermission} & Id, Role, PermissionId \\
    \hline
    \multirow{9}{2.8cm}{\centering Gestion \\ Commerciale} & \textbf{Lead} & Id, CompanyName, ContactName, Phone, Email, Status, AgentId \\
    \cline{2-3}
    & \textbf{Call} & Id, AgentId, LeadId, Transcription, ScorePercentage, Sentiment, CallDuration \\
    \cline{2-3}
    & \textbf{CrmAppointment} & Id, AgentId, ClientName, AppointmentDate, QualityScore, Status \\
    \cline{2-3}
    & \textbf{Appointment} & Id, CallId, DetectedDate, ConfidenceScore, Status \\
    \cline{2-3}
    & \textbf{Campaign} & Id, Name, CompanyName \\
    \cline{2-3}
    & \textbf{Followup} & Id, AgentName, AppointmentDate, Status, RelanceCount \\
    \cline{2-3}
    & \textbf{LeadFolder} & Id, Name, Campaign \\
    \cline{2-3}
    & \textbf{Qualification} & Id, Code, Label, ExpectedKeywords \\
    \cline{2-3}
    & \textbf{Supplier} & Id, \textit{(données fournisseur)} \\
    \hline
    \multirow{3}{2.8cm}{\centering Présence \\ \& Paie} & \textbf{Attendance} & Id, UserId, Date, ClockIn, ClockOut?, Status \\
    \cline{2-3}
    & \textbf{AttendanceBreak} & Id, AttendanceId, BreakStart, BreakEnd?, Type \\
    \cline{2-3}
    & \textbf{Salary} & Id, AgentId, Month, BaseSalary, TotalSalary, PaymentStatus \\
    \cline{2-3}
    & \textbf{SalaryRule} & Id, RuleName, RuleType, Amount, Role, IsActive \\
    \hline
    \multirow{4}{2.8cm}{\centering Qualité \\ \& IA} & \textbf{ManualEvaluation} & Id, AgentId, EvaluatorId, GlobalScore, Decision, Commentaires \\
    \cline{2-3}
    & \textbf{AlertRule} & Id, RuleType, ThresholdValue, IsActive \\
    \cline{2-3}
    & \textbf{AlertHistory} & Id, AgentName, AlertType, Severity, ActualValue \\
    \cline{2-3}
    & \textbf{AiEligibilityLog} & Id, AgentId, ClientData, Result \\
    \hline
    \multirow{3}{2.8cm}{\centering Communication \\ \& Système} & \textbf{Message} & Id, SenderId, ReceiverId, Content, IsUrgent, IsRead \\
    \cline{2-3}
    & \textbf{AgentSavedData} & Id, AgentId, \textit{(données sauvegardées)} \\
    \cline{2-3}
    & \textbf{Log / SourceFile} & \textit{(logs et fichiers)} \\
    \hline
    \end{tabular}

    \vspace{0.3cm}
    \textbf{Relations principales~:}\\[0.1cm]
    \begin{tabular}{@{}ll@{}}
    \hline
    \textbf{Association} & \textbf{Multiplicité} \\
    \hline
    User $\longleftrightarrow$ Attendance & 1 $\longrightarrow$ * \\
    Attendance $\longleftrightarrow$ AttendanceBreak & 1 $\longrightarrow$ * \\
    User $\longleftrightarrow$ CrmAppointment & 1 $\longrightarrow$ * \\
    User $\longleftrightarrow$ Salary & 1 $\longrightarrow$ * \\
    User $\longleftrightarrow$ ManualEvaluation (Agent) & 1 $\longrightarrow$ * \\
    User $\longleftrightarrow$ ManualEvaluation (Evaluator) & 1 $\longrightarrow$ * \\
    User $\longleftrightarrow$ Message (Sender/Receiver) & 1 $\longrightarrow$ * \\
    Lead $\longleftrightarrow$ User & * $\longrightarrow$ 1 \\
    Call $\longleftrightarrow$ Lead & * $\longrightarrow$ 1 \\
    Call $\longleftrightarrow$ Appointment & 1 $\longrightarrow$ * \\
    Permission $\longleftrightarrow$ RolePermission & 1 $\longrightarrow$ * \\
    \hline
    \end{tabular}
    \vspace{0.3cm}
    }}
    \caption{Diagramme de classes global de la plateforme CRM}
    \label{fig:class_diagram}
\end{figure}

Les 26 entités sont réparties en cinq paquetages~: \textbf{Authentification et Sécurité}
(User, Permission, RolePermission), \textbf{Gestion Commerciale} (Lead, Call, CrmAppointment,
Appointment, Campaign, Followup, LeadFolder, Qualification, Supplier),
\textbf{Présence et Paie} (Attendance, AttendanceBreak, Salary, SalaryRule),
\textbf{Qualité et IA} (ManualEvaluation, AlertRule, AlertHistory, AiEligibilityLog, Agent,
AgentSavedData), et \textbf{Communication} (Message, SourceFile, Log).

%==========================================================
\section{Product Backlog}
\label{sec:product_backlog}
%==========================================================

Le Product Backlog est organisé selon la méthodologie Scrum. Les fonctionnalités sont décrites
sous forme d'\textbf{Epics} (objectifs métier) et de \textbf{User Stories} (besoins utilisateurs).
\cite{ScrumGuide}

\begingroup
\scriptsize
\setlength{\tabcolsep}{2pt}
\renewcommand{\arraystretch}{1.18}
\setlength{\LTpre}{6pt}
\setlength{\LTpost}{6pt}

\begin{longtable}{|>{\raggedright\arraybackslash}p{2.7cm}
                   |>{\centering\arraybackslash}p{1.1cm}
                   |>{\raggedright\arraybackslash}p{7.0cm}
                   |>{\centering\arraybackslash}p{1.25cm}
                   |>{\centering\arraybackslash}p{1.35cm}|}
\caption{Product Backlog}\label{tab:backlog}\\
\hline
\textbf{Epic} & \textbf{US ID} & \textbf{User Story} & \textbf{Priorité} & \textbf{Durée estimée (jours)} \\
\hline
\endfirsthead
\hline
\textbf{Epic} & \textbf{US ID} & \textbf{User Story} & \textbf{Priorité} & \textbf{Durée estimée (jours)} \\
\hline
\endhead
\hline
\endfoot
\hline
\endlastfoot
\multirow{8}{2.7cm}{Authentification et gestion des comptes} & 1.1 & En tant qu'\textbf{utilisateur}, je veux m'authentifier avec JWT afin d'accéder à mon espace de travail. & Élevée & 2 \\
\cline{2-5}
 & 1.2 & En tant qu'\textbf{utilisateur}, je veux fermer ma session afin de sécuriser mon accès. & Moyenne & 0.5 \\
\cline{2-5}
 & 1.3 & En tant que \textbf{Super Admin}, je veux gérer les utilisateurs (CRUD) afin d'administrer les comptes. & Élevée & 3 \\
\cline{2-5}
 & 1.4 & En tant que \textbf{Super Admin}, je veux attribuer des rôles et permissions aux utilisateurs afin de contrôler leurs accès. & Élevée & 2 \\
\cline{2-5}
 & 1.5 & En tant qu'\textbf{Admin RH}, je veux créer des comptes agents afin d'intégrer de nouveaux collaborateurs. & Élevée & 1.5 \\
\cline{2-5}
 & 1.6 & En tant que \textbf{Super Admin}, je veux configurer les permissions granulaires afin de définir les droits d'accès. & Moyenne & 2 \\
\cline{2-5}
 & 1.7 & En tant qu'\textbf{Admin RH}, je veux désactiver un compte agent afin de suspendre son accès. & Moyenne & 1 \\
\cline{2-5}
 & 1.8 & En tant qu'\textbf{utilisateur}, je veux réinitialiser mon mot de passe afin de retrouver l'accès à mon compte. & Moyenne & 1 \\
\hline
\multirow{10}{2.7cm}{Gestion des leads} & 2.1 & En tant qu'\textbf{Agent}, je veux créer un lead afin d'ajouter un nouveau prospect. & Élevée & 2 \\
\cline{2-5}
 & 2.2 & En tant qu'\textbf{Agent}, je veux consulter la liste de mes leads afin de suivre mes prospects. & Élevée & 1.5 \\
\cline{2-5}
 & 2.3 & En tant qu'\textbf{Agent}, je veux modifier un lead afin de mettre à jour ses informations. & Élevée & 1 \\
\cline{2-5}
 & 2.4 & En tant qu'\textbf{Agent}, je veux qualifier un lead afin de suivre son état d'avancement. & Élevée & 1.5 \\
\cline{2-5}
 & 2.5 & En tant qu'\textbf{Agent}, je veux filtrer les leads par statut et campagne afin de prioriser mon travail. & Moyenne & 1 \\
\cline{2-5}
 & 2.6 & En tant que \textbf{Supervisor}, je veux affecter un lead à un agent afin de distribuer la charge de travail. & Élevée & 1.5 \\
\cline{2-5}
 & 2.7 & En tant qu'\textbf{Agent}, je veux organiser mes leads dans des folders afin de les classer. & Moyenne & 1 \\
\cline{2-5}
 & 2.8 & En tant que \textbf{Supervisor}, je veux consulter les statistiques de conversion des leads afin d'analyser les performances. & Élevée & 2 \\
\cline{2-5}
 & 2.9 & En tant qu'\textbf{Agent}, je veux planifier un follow-up sur un lead afin d'assurer un suivi régulier. & Moyenne & 1.5 \\
\cline{2-5}
 & 2.10 & En tant qu'\textbf{Agent}, je veux consulter l'historique des interactions d'un lead afin d'avoir une vue complète. & Moyenne & 1 \\
\hline
\multirow{5}{2.7cm}{Rendez-vous et appels} & 3.1 & En tant qu'\textbf{Agent}, je veux créer un rendez-vous pour un lead afin de planifier une rencontre. & Élevée & 2 \\
\cline{2-5}
 & 3.2 & En tant qu'\textbf{Agent}, je veux consulter mon agenda afin de visualiser mes rendez-vous. & Élevée & 1.5 \\
\cline{2-5}
 & 3.3 & En tant qu'\textbf{Agent}, je veux enregistrer un appel afin de tracer l'interaction. & Élevée & 2 \\
\cline{2-5}
 & 3.4 & En tant qu'\textbf{Agent}, je veux consulter l'historique de mes appels afin de suivre mon activité. & Moyenne & 1 \\
\cline{2-5}
 & 3.5 & En tant que \textbf{Quality Manager}, je veux écouter un appel enregistré afin de l'évaluer. & Élevée & 2 \\
\hline
\multirow{8}{2.7cm}{IA et Scoring} & 4.1 & En tant que \textbf{système}, je veux transcrire automatiquement les appels via Whisper afin de produire un texte exploitable. & Élevée & 4 \\
\cline{2-5}
 & 4.2 & En tant que \textbf{système}, je veux scorer automatiquement chaque appel via l'IA afin d'évaluer sa qualité. & Élevée & 4 \\
\cline{2-5}
 & 4.3 & En tant qu'\textbf{Agent}, je veux consulter mon score IA par appel afin de connaître ma performance. & Élevée & 1 \\
\cline{2-5}
 & 4.4 & En tant que \textbf{Quality Manager}, je veux visualiser les scores IA détaillés par critère afin d'analyser la qualité. & Élevée & 1.5 \\
\cline{2-5}
 & 4.5 & En tant que \textbf{Super Admin}, je veux configurer les poids des critères de scoring afin d'adapter l'évaluation. & Élevée & 1 \\
\cline{2-5}
 & 4.6 & En tant que \textbf{système}, je veux analyser le sentiment des appels afin de détecter les tendances. & Moyenne & 3 \\
\cline{2-5}
 & 4.7 & En tant que \textbf{Supervisor}, je veux consulter un dashboard IA agrégé afin de visualiser les performances globales. & Élevée & 2 \\
\cline{2-5}
 & 4.8 & En tant que \textbf{Quality Manager}, je veux comparer le score IA au score manuel afin de valider la pertinence. & Moyenne & 1.5 \\
\hline
\multirow{6}{2.7cm}{Qualité et évaluations} & 5.1 & En tant que \textbf{Quality Manager}, je veux créer une évaluation manuelle d'un appel afin de noter la qualité. & Élevée & 2 \\
\cline{2-5}
 & 5.2 & En tant que \textbf{Quality Manager}, je veux consulter l'historique des évaluations afin de suivre les progrès. & Moyenne & 1 \\
\cline{2-5}
 & 5.3 & En tant que \textbf{Super Admin}, je veux configurer des règles d'alerte qualité afin d'être notifié des dérives. & Élevée & 1.5 \\
\cline{2-5}
 & 5.4 & En tant que \textbf{Quality Manager}, je veux recevoir une alerte quand un score passe sous un seuil afin d'intervenir rapidement. & Moyenne & 1 \\
\cline{2-5}
 & 5.5 & En tant que \textbf{Supervisor}, je veux consulter un dashboard qualité afin de piloter la performance de l'équipe. & Élevée & 2 \\
\cline{2-5}
 & 5.6 & En tant que \textbf{Quality Manager}, je veux exporter les résultats d'évaluation afin de les partager. & Moyenne & 1 \\
\hline
\multirow{5}{2.7cm}{Présence et paie} & 6.1 & En tant qu'\textbf{Agent}, je veux pointer mon arrivée et mon départ afin de déclarer ma présence. & Élevée & 2 \\
\cline{2-5}
 & 6.2 & En tant qu'\textbf{Agent}, je veux déclarer une pause afin de tracer mes absences. & Moyenne & 1 \\
\cline{2-5}
 & 6.3 & En tant qu'\textbf{Admin RH}, je veux consulter les pointages de présence afin de vérifier les horaires. & Élevée & 1.5 \\
\cline{2-5}
 & 6.4 & En tant qu'\textbf{Admin RH}, je veux configurer les règles de calcul des salaires afin d'automatiser la paie. & Élevée & 2.5 \\
\cline{2-5}
 & 6.5 & En tant qu'\textbf{Admin RH}, je veux générer les bulletins de paie afin de préparer la rémunération. & Élevée & 2 \\
\hline
\multirow{4}{2.7cm}{Campagnes et analytics} & 7.1 & En tant que \textbf{Supervisor}, je veux créer une campagne afin d'organiser une action commerciale. & Élevée & 2 \\
\cline{2-5}
 & 7.2 & En tant que \textbf{Supervisor}, je veux suivre les performances d'une campagne afin d'évaluer son efficacité. & Élevée & 2 \\
\cline{2-5}
 & 7.3 & En tant que \textbf{Supervisor}, je veux consulter les KPIs du centre d'appels afin de piloter l'activité. & Élevée & 2 \\
\cline{2-5}
 & 7.4 & En tant que \textbf{Supervisor}, je veux exporter les rapports en PDF afin de les diffuser. & Moyenne & 1.5 \\
\hline
\end{longtable}
\endgroup

%==========================================================
\section{Répartition des sprints}
\label{sec:repartition_sprints}
%==========================================================

Le tableau~\ref{tab:sprints} résume la planification selon une progression en trois releases.
La durée totale du stage est structurée en \textbf{6 sprints de 20 jours}, soit \textbf{120 jours}.

\begin{table}[H]
    \centering
    \caption{Planification des sprints sur trois releases}
    \label{tab:sprints}
    \scriptsize
    \setlength{\tabcolsep}{3pt}
    \renewcommand{\arraystretch}{1.35}
    \begin{tabular}{|>{\raggedright\arraybackslash}m{2.8cm}
                     |>{\centering\arraybackslash}m{1.6cm}
                     |>{\raggedright\arraybackslash}m{8.3cm}
                     |>{\centering\arraybackslash}m{2.2cm}|}
        \hline
        \rowcolor[gray]{0.80}
        \textbf{Phase / Release} & \textbf{Sprint} & \textbf{Fonctionnalités du Sprint} & \textbf{Estimation en jours} \\
        \hline
        \multirow{2}{2.65cm}{\textbf{Release 1 : Socle de la plateforme}} & Sprint 1 &
        \begin{itemize}[leftmargin=*,nosep]
            \item Authentification JWT, gestion des utilisateurs et des rôles.
            \item Permissions granulaires et RBAC.
        \end{itemize} & 20 jours \\
        \cline{2-4}
        & Sprint 2 &
        \begin{itemize}[leftmargin=*,nosep]
            \item Gestion des leads (CRUD, qualification, folders, follow-ups).
            \item Rendez-vous, appels et campagnes.
        \end{itemize} & 20 jours \\
        \hline
        \multirow{2}{2.65cm}{\textbf{Release 2 : IA et Qualité}} & Sprint 3 &
        \begin{itemize}[leftmargin=*,nosep]
            \item Intégration Whisper pour la transcription des appels.
            \item Moteur de scoring IA basé sur gemma3:4b avec critères pondérés.
        \end{itemize} & 20 jours \\
        \cline{2-4}
        & Sprint 4 &
        \begin{itemize}[leftmargin=*,nosep]
            \item Évaluation qualité manuelle et automatique.
            \item Alertes, règles qualité et dashboard qualité.
        \end{itemize} & 20 jours \\
        \hline
        \multirow{2}{2.65cm}{\textbf{Release 3 : Consolidation}} & Sprint 5 &
        \begin{itemize}[leftmargin=*,nosep]
            \item Gestion de présence (pointage, pauses).
            \item Calcul automatique des salaires.
        \end{itemize} & 20 jours \\
        \cline{2-4}
        & Sprint 6 &
        \begin{itemize}[leftmargin=*,nosep]
            \item Tableaux de bord analytics et KPIs.
            \item Exports PDF, CI/CD et déploiement Docker.
        \end{itemize} & 20 jours \\
        \hline
    \end{tabular}
\end{table}

%----------------------------------------------------------
\section{Conclusion}
\phantomsection
\label{sec:chap2_conclusion}
%----------------------------------------------------------

Ce chapitre a présenté l'analyse complète des besoins de la plateforme CRM intelligente. Les
cinq acteurs (\textbf{Agent}, \textbf{Quality Manager}, \textbf{Admin RH}, \textbf{Supervisor}
et \textbf{Super Admin}) ont été identifiés avec une description détaillée des cas
d'utilisation par rôle. Le Product Backlog, organisé en \textbf{7 Epics métier} et \textbf{6
sprints} planifiés sur trois releases, constitue la feuille de route du projet. Le chapitre
suivant détaille l'environnement de travail et les choix technologiques retenus.
